\chapter{Discussion}
\label{sec:discussion}

The previous chapter ended on a gap, reward models that look strong on every standard offline metric and still fail the moment a policy optimizes against them. This chapter places that finding, and our results more broadly, against the prior work, then takes an honest inventory of what limits them and what they imply beyond this thesis.

\section{Reading the results against prior work}

Our results do not contradict the published claims of the models we compare against. Robometer reports strong offline reward quality and real-world RL with its released model, and our evaluation reproduces its offline numbers on the held-out-robot set, which is what lets us trust the rest of our harness. What Robometer never reports is the experiment at the center of this thesis, the released model driving simulated RL without in-domain adaptation, and in that experiment it fails on almost every task we test. The same holds down the table. RoboReward was only ever validated in the real world, RoboDopamine adapts to each task from a demonstration before learning, and LRM, the one prior model tested zero-shot in simulated RL, pays our policies the richest reward in the comparison while producing almost no success. Each system behaves as its paper says it does. The problem is that the evaluations those papers chose did not surface the failure mode that dominates once a policy is in the loop.

One could ask whether our results reflect a plain real-to-sim gap, models that simply cannot read simulated frames. We doubt that this is the main source. These models already train on simulated data, Robometer's corpus includes the LIBERO suites and MetaWorld clips among its sources, and the released model's perfect Close-Drawer run shows it can score simulated frames well enough to train a policy to full success when the regime suits it. A rendering gap may contribute at the margin, but the failure mode we observe, reward paid rising while true success stays flat, is exploitation, not blindness. The simulated setting also has a virtue beyond ground truth. Real-world experiments are hard to reproduce, while simulated ones can be rerun exactly, so we would encourage future reward models to be tested thoroughly in simulation as well, where such claims can be checked by anyone.

That failure mode is not new in kind. Learned reward functions have been known to invite exploitation since the early safety literature \citep{amodei_concrete_2016}, and reinforcement learning from human feedback showed that optimizing a policy against a learned proxy eventually degrades true performance even as the proxy climbs \citep{gao_scaling_2023}. What our experiments add is a concrete robotic instance in which the mechanism is directly observable, because simulation provides ground-truth success independent of the reward. The paired curves of Figure~\ref{fig:mw-reward}, reward paid rising while true success stays at zero, are overoptimization made visible, and we would argue that reporting this pair should become standard practice for reward-model papers, since either curve alone tells a misleading story.

Our asymmetric objective, in turn, is the training-time counterpart of an observation made at inference time by \citet{gumbsch_demonstrations_2026}, that false positives mislead an optimizing policy far more than false negatives. Building that asymmetry into the loss is what separates the surviving fine-tunes from the collapsing ones in both simulators, and it costs almost nothing to adopt, a reweighting of an existing objective that resumes directly from pretrained weights.

It is worth spelling out why these two ingredients help, because the mechanism carries beyond our models. A corpus of successful demonstrations teaches a reward what completed tasks look like, but a policy in training spends nearly all of its time in states no demonstration contains, half-finished reaches, near-misses, and outright failures, and on those states a success-trained model is free to guess, usually upward. Our failure data fills exactly the region the policy actually visits, which is why the data alone lifts a reward from never moving a policy on Robomimic to training one. The asymmetric loss adds the complementary piece. The most dangerous residual errors are the near-misses that look almost complete, and pricing false positives above false negatives pushes precisely those below the firing threshold. The data covers the states an optimizing policy reaches, the loss decides how to err on them, and it is the combination that survives optimization.

The dense progress labels have to pass the same test. Offline they deliver, and the ablation of Section~\ref{sec:results-icl} ties those gains to the data they came from, but as a live reward the progress head never made its case, a shortfall we return to with the limitations below.

\section{Why the offline metrics missed it}

Offline evaluation scores a reward on trajectories drawn from a fixed distribution. A policy under training is not such a distribution. It is a search process that shifts probability toward whatever the reward scores highly, so a region of false positives that occupies a negligible slice of a test set becomes, once discovered, most of what the reward is asked to judge. The shift is not even neutral. For overestimated states reinforcement learning is adversarial to the reward model, since every error the policy discovers is revisited, trained on, and exploited harder, a feedback loop that self-amplifies exactly the mistakes a fixed test set touches only once. Metrics that integrate over the whole score range, ROC-AUC above all, are close to blind to this, because they average over operating points a deployed detector will never use. The quantities that came closest to predicting our downstream results were the ones tied to a conservative operating point, the true-positive rate at a strict false-positive budget and the class separation $d'$, and even these were only necessary signals, not sufficient ones. RoboRef-Std posted a respectable $d'$ of $1.08$ and still came apart under an optimizing policy. We looked for an offline number that cleanly forecasts on-policy survival and did not find one, which is itself a result: with the metrics the field currently reports, the downstream test cannot be skipped.

\section{Limitations}

The claim we can defend is direction, not finality. We identified the dominant failure mode of current vision-language reward models, built a dataset, an input format, and an objective that target it, and the resulting models beat every published reward model we tested in most of the RL regimes we evaluate. What we cannot claim is stability. A reward model's behavior shifts with the checkpoint and the threshold it is read at, off-policy actor-critic methods are notoriously sensitive on their own, and the two sensitivities compound, so individual runs can and do surprise, as the early rise and collapse of two of our own fine-tunes on Robomimic shows. This bears on the reliability of any single curve we report, and it is compounded by seed counts that are honest for a compute-bound academic project but modest by RL standards, five per model on MetaWorld and LIBERO and a single seed on Robomimic, where one run takes several days.

The second limitation is scope, and for a robotics thesis it is the one that matters most: every policy we train lives in simulation. Simulation is what makes the ground-truth success measurements behind our central finding possible at all, but a real robot adds shifting lighting, changing scenes, sensor noise, and no simulator state to calibrate against, and our only evidence about real-world behavior is offline, the held-out-robot evaluation set. The generalizability of the downstream results to physical hardware is therefore untested, and claiming otherwise would repeat exactly the evaluation gap we criticize in prior work.

A further limitation sits inside the dataset contribution itself. The per-frame progress labels are the most expensive part of our corpus, and while the offline chapters justify them, the configuration that trains the best policy reads the success head as a sparse detector in every benchmark. MetaWorld and Robomimic therefore report no progress-head runs, and in the LIBERO regime study, the one place the progress head is deployed as the reward, our models improve on the baseline yet do not train the tasks to a strong success rate. The evidence that the dense labels help downstream is indirect, running through the success head that was trained alongside them, and a policy driven by the progress signal itself is something this thesis does not validate.

Our deployment recipe is also not fully zero-shot. The success threshold is calibrated per model on a small set of rollouts, and although we found the working range to be narrow and stable, between $0.75$ and $0.90$ across simulators and tasks, the procedure still needs some calibration. The minimum episode length is the other per-task number, and it mattered for every vision-language reward model in the comparison, not only ours. It is also easy to set. The length of a demonstration worked in every simulator, and those demonstrations were already there, since IBRL pretrains its policy on them. Outside the in-context model we take no images and no actions from them, only a single number. One can still object that needing that number makes the method less than strictly zero-shot, and the objection is fair, though the same could be said of the maximum episode length that every policy setup fixes per task without anyone counting it as supervision. The next natural step is to move the threshold calibration online, adjusting it from the policy's own behavior during training, and to search for general minimum-length values that hold across a whole suite instead of per task.

\section{Ethical considerations}

Three concerns are worth naming. The first is safety. A reward model that fires falsely on a real robot does not just corrupt a learning curve, it reinforces physical behavior that never achieved the task, and an optimizing policy will repeat whatever was reinforced. Conservatism is therefore not only a performance choice but the safety-aligned default for any learned reward that touches hardware, and we see our results as evidence for making it standard. The second concern is inherited bias. These reward models judge the physical world through backbones pretrained on internet data, and whatever blind spots that data carries, scenes, objects, and environments it underrepresents, propagate silently into which robot behaviors get rewarded. The third is oversight. Our labeling pipeline replaces human annotation with large models precisely because the corpus is too large to label by hand, which also means no human ever verifies most of the labels the reward model learns from. We audited samples and measured error rates, but automated curation at this scale trades away a layer of human review, and that trade should be stated, not hidden. Alongside these concerns sits the compute cost of the whole enterprise, from labeling a million clips with large models to fine-tuning on H100 nodes, an energy footprint that grows with every scaling step the field takes.

The limitations above sketch the road ahead, a real-world reinforcement learning loop as the decisive test, online threshold calibration to close the zero-shot gap, and reward pretraining for newer backbones.
